% Corte mecánico íntegro: parte_iv.tex:409--499.
% Residencia material del ensamblaje sucesor de 102 capítulos.
% Residencia de control tipado: capítulo 90.
\chapter{Invariantes de consistencia matemática y física de las realizaciones helicoidales}

\section{Invariantes matemáticos de la prolongación helicoidal}

La teoría queda refutada en el dominio correspondiente si ocurre cualquiera
de los hechos siguientes:
\begin{enumerate}
\item una vuelta nonádica borra la profundidad en el levantamiento;
\item la regla \(g(K+9)=g(K)\) se confunde con retorno del estado;
\item el lector radial no conserva factores y holonomías prefijas;
\item falla \cref{espiral:eq:naturalidad-truncamiento};
\item la acción de \(\Gamma_9\) no satisface la equivarianza
\cref{espiral:eq:equivarianza-gamma9};
\item el carry compuesto incumple \cref{espiral:eq:cociclo-carry};
\item una curva clásica interviene en la selección de \(p,\Lambda,C\);
\item la ley \cref{espiral:eq:lp-producto} o la inversión
\cref{espiral:eq:lp-inversion} falla en su dominio;
\item el tornillo de Pell no produce \cref{espiral:eq:espiral-log-pell};
\item se identifica \(\tau\), \(p\) o \(\kappa_{\mathrm{F}}\) sin un mapa
tipado.
\end{enumerate}

\section{Condiciones constitutivas de la realización electromagnética y material}

Una realización electromagnética o material queda refutada si:
\begin{enumerate}
\item incumple \(\mathbf E\perp\mathbf B\perp\widehat{\mathbf k}\),
\(|E|=c|B|\) o \(\omega=ck\);
\item interpreta la dilatación de Pell como crecimiento energético sin
fuente;
\item confunde helicidad fotónica con espín fermiónico;
\item no preserva \(J^2=-I\) y \(\star^2=-I\) en el entrelazador declarado;
\item la firma espiral no es natural sobre las multisecciones del atlas;
\item una etiqueta física selecciona retrospectivamente una celda o una
curva;
\item se identifica el eje \(\varphi\) de la genealogía de constantes con
el electrón material sin mapa dimensional.
\end{enumerate}

\section{Estratificación de los teoremas internos y de sus realizaciones físicas}

\begin{longtable}{p{.31\textwidth}p{.22\textwidth}p{.37\textwidth}}
\caption{Estatuto probatorio de los resultados.}
\label{espiral:tab:estatutos-finales}\\
\toprule
Resultado&Estatuto&Alcance\\
\midrule
\endfirsthead
\toprule
Resultado&Estatuto&Alcance\\
\midrule
\endhead
Curvatura mixta APP, regímenes TRIT, \(\Gamma_9\), hélice nonádica y tornillo de Pell&\status{teorema interno previo}&Exacto en los dominios definidos por las construcciones precedentes.\\
Envolvente afín, lector por prefijos, naturalidad, equivarianza e inversión&\status{teorema de esta obra}&Exacto para truncamientos, rutas tipadas y sector reversible; \(\sd_4\) certificado.\\
Familia \(L_p\) y cinco publicaciones&\status{teorema de esta obra}&Clasifica la familia potencia--logarítmica, no todas las curvas.\\
Carta \(p=\tau_++\tau_-\), \(a=\tau_+-\tau_-\)&\status{formalización interna exacta}&Caracteres elementales de la categoría levantada de 2 semicoronas; los cociclos globales se obtienen por suma de ruta.\\
Modo electromagnético circular&\status{realización física}&Solución restringida de Maxwell en el fibrado fase--posición.\\
Firma espiral de partículas&\status{realización condicionada}&Propuesta sobre multisecciones; queda sujeta al cuadrado explícito de constancia, separación y naturalidad material.\\
Electrón central \((5,5)\)&\status{teorema interno previo}&Anclaje material único fijado por la inversión celular.\\
\bottomrule
\end{longtable}

\section{Conclusión}

La teoría obtiene una respuesta unificada a la intuición de partida. Una
órbita es una hoja helicoidal cuando se conserva la profundidad; la
involución bidireccional produce hojas conjugadas de ascenso y descenso; los
círculos anidados son secciones de nivel; los cilindros son familias de esas
secciones. Las espirales clásicas no se eligen como plantillas. Aparecen al
publicar en \(\R^2\) caracteres radiales distintos de una historia
APP--TRIT--TPK.

El resultado principal no es una nueva lista de curvas. Es una separación
estructural entre genealogía y apariencia, acompañada de un lector que
conserva exactamente la información que la apariencia puede olvidar. La
curva se sitúa al final. La recta real, el plano real y las geometrías de
Frenet son publicaciones. La fase puede retornar mientras la memoria avanza;
la circunferencia puede ser intrínseca o proyectada; la espiral logarítmica
puede ser la sombra exacta de un tornillo espectral; y el mismo mecanismo
puede admitir realizaciones arquimedianas, electromagnéticas y materiales sin
confundir sus codominios.

\begin{resultbox}[title={Clasificación functorial del lector radial y estratificación de sus realizaciones}]
La composición APP--TRIT--TPK genera una clase de historias radiales con
holonomía, orientación, carry, frontera y memoria. El lector enriquecido es
natural bajo truncamiento y equivariante bajo la holonomía nonádica. Su
publicación arquimediana produce curvas potencia--logarítmicas y suspensiones
helicoidales, pero no es fiel: una misma figura puede conservar genealogías
distintas. Ésta es la estructura íntima dinámica del espacio que la
clasificación convencional de curvas observa sólo después de la proyección.
\end{resultbox}
